\section{Working in the volatility chart}
\label{sec:svichart}

Chapter~\ref{ch:lqd} built a smile in two steps.  First choose a
probability law for the return; then read prices, and hence implied
volatilities, off that law.  The order of the steps did real work.
Because every curve the model could produce came from a genuine law,
validity never needed checking.  The cost was indirection: the model's
coordinates $(L,R,a_n)$ --- the endpoint speeds and polynomial
coefficients of its rank-space construction --- describe how probability
is transported, and none of them is a number a trading desk quotes.

This chapter works in the opposite order, which is also the industry's
order.  Two definitions from Chapter~\ref{ch:lqd} are used on every page
that follows, so we restate them once.  The \emph{log-moneyness} of a
strike $K$ against the forward $F$ is $k=\log(K/F)$: the money is $k=0$,
and negative $k$ is the put side.  The \emph{total implied variance} at
that strike is $w(k)=\sigma(k)^2\tau$, the squared implied volatility
times the variance time $\tau$.  One expiry's smile is then a single
function $w(k)$, and the direct route is to write that function down as
an explicit formula with a few parameters.  Working this way --- on the
picture whose axes are strike and variance, rather than behind it on the
law --- is what we call working \emph{in the volatility chart}.

The direct route has a real cost, and it is exactly the one
\cref{sec:intro} identified when it introduced the function $g_{\rm D}$.
Most functions $w(k)$ are not the smile of any probability law.  A model
that writes $w(k)$ freely therefore moves through a space that contains
invalid curves, and validity --- which Chapter~\ref{ch:lqd} got for free
--- must now be policed.  Part of the policing can be done by cheap
inequalities that the optimizer respects while it runs.  Part of it
cannot, and needs a check on the finished curve.  Learning which part is
which is most of this chapter.

Why pay that cost?  Because the direct route serves two audiences at
once.  An optimizer wants a formula that evaluates in one line and
differentiates in closed form.  A desk wants parameters it can talk
about: where the at-the-money volatility sits, how steep the skew is,
how heavy each wing is.  The stochastic-volatility-inspired (SVI) family
of Gatheral \citep{Gatheral2004} answers the optimizer with a tilted
hyperbola,
\begin{equation}
\boxed{\;w(k)=a_S+b_S\Big\{\rho_S(k-m_S)+\sqrt{(k-m_S)^2+s_S^2}\Big\},\;}
\label{eq:svi}
\end{equation}
five parameters under the structural conditions $b_S>0$, $|\rho_S|<1$,
$s_S>0$.  (The subscript $S$ is there to keep the book's notation
unambiguous: bare $\rho$ is Chapter~\ref{ch:lqd}'s logistic density and
bare $m$ its martingale shift, so the SVI parameters carry their own
mark.)  The desk is answered by a second coordinate system on the same
curve --- the jump-wing handles of \cref{sec:svijw}, five numbers chosen
to match the desk's questions.  The two coordinate systems describe the
same family, and \cref{thm:jwimage} will say exactly when one converts
to the other.

Before any geometry, we need to be precise about what ``valid'' means
here, because different validity conditions will be supplied by
different mechanisms.  Five graded conditions recur through the chapter;
\cref{tab:tiers} collects them once, and the sections that follow fill
the table in.  The first three are inequalities on the parameters, cheap
enough to build into the fit itself.  The fourth is the full butterfly
condition $g_{\rm D}\ge0$ of \eqref{eq:durrleman} --- recall that
$g_{\rm D}(k)$ is, up to a positive factor, the probability density the
smile implies at strike $k$, so a valid smile needs $g_{\rm D}\ge0$
everywhere.  No cheap inequality implies it: \cref{sec:svibelly}
exhibits a slice that passes every screen and still fails it.  The fifth
condition couples different expiries and is not new to this chapter: it
is the calendar coupling of \cref{sec:calendar}, which applies to these
families unchanged.

The table names two symbols and two mechanisms that are still ahead of
us, so read it as a map of the chapter rather than a summary of things
known: $\beta_L$ and $\beta_R$ are the smile's two wing slopes, defined
in \cref{sec:svigeom}; a \emph{hinge penalty} is a cheap fence built
into the fit, explained in \cref{sec:svibelly}; and the
\emph{structural chart} is a choice of optimizer coordinates, built in
\cref{sec:svicalib}.  We will return to this table as each row is
earned.

\begin{table}[htbp]
\centering\small
\setlength{\tabcolsep}{5pt}
\renewcommand{\arraystretch}{1.15}
\begin{tabular}{@{}p{0.20\textwidth}p{0.40\textwidth}p{0.32\textwidth}@{}}
\toprule
Condition & Mathematical statement & Supplied by\\
\midrule
Structural & $b_S>0$, $|\rho_S|<1$, $s_S>0$ & the optimization chart
(\cref{sec:svicalib})\\
Positive floor & $\min_k w(k)\ge0$ & a hinge penalty; automatic under the
structural chart\\
Wing-admissible & $\max(\beta_L,\beta_R)\le\beta_{\max}<2$ & a hinge
penalty; automatic under the structural chart\\
Butterfly-free & $g_{\rm D}(k)\ge0$ for all $k$ & \emph{not} implied by
the rows above; checked on a dense grid (\cref{sec:svibelly})\\
Calendar-ordered & $w_{\tau_1}(k)\le w_{\tau_2}(k)$, $\tau_1<\tau_2$ &
the coupling of \cref{sec:calendar}, applied model-agnostically\\
\bottomrule
\end{tabular}
\caption{Five grades of validity for a slice written directly in the
volatility chart.  In Chapter~\ref{ch:lqd} the first four held by
construction.  Here each must be attributed to the mechanism that
supplies it --- and the fourth row is the one no cheap mechanism
reaches.}
\label{tab:tiers}
\end{table}

The chapter now proceeds in order.  \Cref{sec:svigeom} works out the
geometry of the hyperbola \eqref{eq:svi}: its vertex, its wings, and a
rigidity that will matter later.  \Cref{sec:svilee} explains why the
wing slopes must be capped, and why the cap sits strictly inside Lee's
bound.  \Cref{sec:svibelly} shows what the caps do \emph{not} protect
--- the middle of the smile --- and states the chapter's division of
labor between fences and checks.  \Cref{sec:svijw} develops the trader
coordinates and proves exactly where they can be inverted.
\Cref{sec:svicalib} turns to calibration: what the least-squares
objective integrates against, and how to choose coordinates so the
cheap conditions hold at every step.  \Cref{sec:mcs} then builds the
Multi-Core Sigmoid, which adds capacity in the body of the smile without
touching the wings.  \Cref{sec:comparison} closes by fitting LQD, SVI,
and MCS to the same frozen quotes under one protocol.
